% =============================================================================
%  Weeks 5-7, step by step.
%  A narrative walkthrough written to be read aloud in a supervision meeting.
% =============================================================================
\documentclass[11pt,a4paper]{article}

\usepackage[T1]{fontenc}
\usepackage[utf8]{inputenc}
\usepackage[margin=2.3cm]{geometry}
\usepackage{booktabs}
\usepackage{array}
\usepackage{xcolor}
\usepackage{graphicx}
\usepackage{enumitem}
\usepackage{caption}
\usepackage{amsmath}
\usepackage{amssymb}
\usepackage{float}
\usepackage{titlesec}
\usepackage[hidelinks]{hyperref}

\graphicspath{{../outputs/model/figures/}{../outputs/hbayes/figures/}}

\definecolor{stepblue}{HTML}{1F4E79}
\definecolor{boxbg}{HTML}{F4F6F8}
\definecolor{askgrey}{HTML}{52514E}

\newcommand{\code}[1]{\texttt{\small #1}}
\newcommand{\arr}{$\rightarrow$}

\setlist[itemize]{itemsep=2pt,topsep=3pt}
\setlist[enumerate]{itemsep=2pt,topsep=3pt}
\captionsetup{font=small,labelfont=bf}

\titleformat{\section}{\color{stepblue}\Large\bfseries}{\thesection.}{0.6em}{}
\titlespacing*{\section}{0pt}{18pt}{8pt}

% "What we found" box
\newenvironment{found}
  {\begin{center}\begin{minipage}{0.95\textwidth}%
   \vspace{4pt}\hrule\vspace{6pt}\small}
  {\vspace{6pt}\hrule\vspace{4pt}\end{minipage}\end{center}}

% "If the professor asks" block
\newcommand{\ifasked}[2]{%
  \vspace{4pt}\noindent{\small\color{askgrey}\textbf{If asked: ``#1''}\\ #2}\vspace{3pt}}

\title{\textbf{What We Did in Weeks 5--7}\\[0.5em]
       \large A step-by-step walkthrough}
\author{Parv Patel \quad\textbullet\quad Mayank Tewatia\\[0.3em]
        \small Department of Data Science, IIT Palakkad}
\date{August 2026}

\begin{document}
\maketitle

\vspace{-1em}
\begin{found}
\textbf{The one-sentence summary, if you only get to say one thing.}

We built the Markov model the plan asked for, found it did not work, tested
\emph{why} rather than guessing, and the test told us the observed subtype is a
noisy reading of a hidden state. The model built on that finding beats the
baseline in 10 out of 10 validation splits.
\end{found}

\noindent
This document walks through the fortnight in the order the work actually
happened, so it can be read straight through. Each step has three parts: what
we did, why we did it, and what we found. Where a question is likely, the
answer is given underneath.

\clearpage
\tableofcontents
\clearpage

% =============================================================================
\section{Where we were starting from}
% =============================================================================

\textbf{What we had.} A dataset of 1{,}080 Parkinson's patients, 4{,}441 visits
where we could compute the TD/PIGD subtype, and 3{,}361 observed transitions
between subtypes. Every label had been checked against PPMI's own
classification and agreed 100.00\% on the two-class version and 99.87\% on the
three-class version.

\textbf{What we already knew from the previous fortnight.} Three things, and all
three turned out to matter:

\begin{enumerate}
  \item Indeterminate is barely a state --- only 19.2\% of patients are still
        Indeterminate at their next visit.
  \item 35.2\% of all observed subtype changes reverse at the very next visit.
  \item The chance of changing subtype depends strongly on how long the gap
        between visits was.
\end{enumerate}

\textbf{What we were asked to do.} Build a simple probabilistic model --- a
Markov model --- that says, for a patient in a given subtype, how likely they
are to stay or move. Then validate it on patients the model has never seen.

% =============================================================================
\section{Step 1 --- We built the simple model first}
% =============================================================================

\textbf{What we did.} A discrete Markov chain. This is the model the plan
describes. It is one table of probabilities: given the subtype at this visit,
the probability of each subtype at the next visit.

\textbf{Why.} It is the simplest thing that answers the question, and it is what
we were asked for. Starting anywhere more complicated would mean we could not
say whether the complexity was earning anything.

\textbf{What we found.} It works reasonably. On patients held out from fitting,
it predicts the next subtype correctly 76.8\% of the time, and comfortably beats
a null model that ignores the current subtype entirely.

\ifasked{Why start with something so basic?}{Because if we had started with a
complicated model and it worked, we would not know whether the complication was
necessary. Every model after this one had to beat it on data it had never seen.}

% =============================================================================
\section{Step 2 --- We built the model we expected to be better}
% =============================================================================

\textbf{What we did.} A continuous-time Markov chain. Instead of one table, this
estimates a matrix of \emph{rates}, and gives the transition probability for
\emph{any} elapsed time through $P(t) = e^{Qt}$.

\textbf{Why.} Two reasons. First, the previous fortnight showed the chance of
changing subtype depends heavily on the gap between visits --- so a model that
ignores the gap is throwing away information. Second, a clinician needs answers
at 1, 2, 3 and 5 years, and a per-visit table cannot give those, because visits
are not one year apart.

\textbf{What we found.} \textbf{It was worse.} Not slightly --- it lost at every
interval length, and it lost by the most at long intervals, which is exactly
where accounting for the gap should have helped the most.

\begin{table}[H]
\centering
\begin{tabular}{@{}lrrr@{}}
\toprule
\textbf{Gap between visits} & \textbf{$n$} & \textbf{Continuous-time} & \textbf{Discrete} \\
\midrule
Up to 1.5 years & 890 & $-0.6534$ & $-0.6488$ \\
1.5 to 2.5 years & \phantom{0}71 & $-0.8151$ & $-0.7907$ \\
Over 2.5 years & \phantom{0}41 & $-0.9054$ & $-0.8061$ \\
\bottomrule
\end{tabular}
\caption{Held-out score, higher is better. The continuous-time model is behind
everywhere, and furthest behind on long gaps.}
\end{table}

\ifasked{So the previous fortnight's finding was wrong?}{No. The finding was
that transition probability depends on the gap, and that is still true. What
this shows is that \emph{correcting} for the gap inside this particular model
does not help --- which means something more basic is wrong with the model.
That is what we investigated next.}

% =============================================================================
\section{Step 3 --- Instead of guessing, we tested the assumption}
% =============================================================================

\textbf{What we did.} Both models rest on the same assumption, called the Markov
property: that knowing a patient's subtype \emph{today} tells you everything
useful --- where they were before should not matter.

That is a statement we can check directly. We took every patient with three
labelled visits in a row --- 2{,}427 of them --- and asked: given where they are
now, does where they came from change what happens next?

\textbf{Why.} Because if we had simply tried a third model and it worked better,
we would not know why. A failed model is much more useful if you can say what
broke.

\textbf{What we found.} The assumption is false, and not marginally.

\begin{table}[H]
\centering
\begin{tabular}{@{}lrrr@{}}
\toprule
\textbf{Patients currently labelled TD, who arrived from\ldots} & \textbf{\arr\ TD} & \textbf{\arr\ Indet.} & \textbf{\arr\ PIGD} \\
\midrule
TD (they were TD before too)   & 84.9\% & 6.5\% & \textbf{8.6\%} \\
PIGD (they were PIGD before)   & 48.2\% & 12.5\% & \textbf{39.3\%} \\
\bottomrule
\end{tabular}
\caption{$\chi^2 = 156.6$, $p = 7.9 \times 10^{-33}$. Rejected in all three states.}
\end{table}

\begin{found}
\textbf{Read this row out loud if nothing else.} Two patients are both labelled
TD today. One was TD last time; the other was PIGD last time. The first has an
8.6\% chance of being PIGD next visit. The second has a 39.3\% chance.

\textbf{That is a 4.6-fold difference, and the Markov assumption says it should
be zero.}
\end{found}

\textbf{What that pattern means.} If a patient's true state is PIGD but one
examination happens to record TD, the next examination will probably record PIGD
again. The patient's history is carrying information that the current
\emph{label} has lost. That is the signature of a hidden state being observed
through a noisy measurement.

\ifasked{Could this just be noise in the test?}{No --- $p = 7.9\times10^{-33}$
on 1{,}583 observations for the TD row alone, and the test is rejected
independently in all three states.}

% =============================================================================
\section{Step 4 --- We built the model the diagnosis pointed to}
% =============================================================================

\textbf{What we did.} A hidden Markov model. It has two layers:

\begin{itemize}
  \item A \textbf{true phenotype} that we never observe directly, which changes
        over time at estimated rates.
  \item A \textbf{measurement} of that phenotype at each visit, which can be
        wrong, with the chance of being wrong also estimated from the data.
\end{itemize}

\textbf{Why.} Because that is exactly the structure the Step 3 test pointed to,
and --- worth saying to the professor --- it is the same structure the project
proposal already specifies, in the line \emph{observed ratio = true ratio +
measurement noise}. We did not invent a new direction; we arrived at the one
already written down, from the data.

\textbf{What we found.} It wins, and not on one lucky split.

\begin{table}[H]
\centering
\begin{tabular}{@{}lrrr@{}}
\toprule
\textbf{Model} & \textbf{Score} & \textbf{Accuracy} & \textbf{Brier} \\
\midrule
Hidden Markov (true state $+$ measurement error) & $\mathbf{-0.6384}$ & \textbf{77.0\%} & \textbf{0.346} \\
Discrete Markov chain & $-0.6653$ & 76.8\% & 0.363 \\
Continuous-time Markov chain & $-0.6752$ & 76.6\% & 0.367 \\
Null model & $-0.8633$ & 62.5\% & 0.519 \\
\bottomrule
\end{tabular}
\end{table}

\begin{found}
We repeated the entire procedure --- new random split of patients, refit every
model from scratch --- \textbf{ten times}. The hidden Markov model won
\textbf{10 out of 10}. Mean advantage $+0.028 \pm 0.009$ per prediction, paired
$t$-test $p = 5.4\times10^{-6}$.
\end{found}

\begin{figure}[H]
\centering
\includegraphics[width=\textwidth]{model_comparison.png}
\caption{Each grey line is one split, joining the four models. No line crosses
between the hidden Markov model and the rest, so the ranking is consistent
rather than an average of mixed results.}
\end{figure}

\ifasked{A more complicated model beating a simpler one --- isn't that just overfitting?}{%
Overfitting shows up as doing better on the data you fitted and worse on data
you did not. Every number here is on patients the model never saw during
fitting, and the win repeats on ten independent splits. If it were overfitting,
it would lose out of sample, not win ten times.}

% =============================================================================
\section{Step 5 --- What the model tells us about the disease}
% =============================================================================

Two results here are worth presenting as findings in their own right.

\subsection*{Progression is sequential --- patients pass through the middle}

\begin{table}[H]
\centering
\begin{tabular}{@{}lrrr@{}}
\toprule
\textbf{From} & \textbf{\arr\ TD} & \textbf{\arr\ Indeterminate} & \textbf{\arr\ PIGD} \\
\midrule
TD            & --- & 0.0937 & \textbf{0.0057} \\
Indeterminate & 0.1156 & --- & 0.3557 \\
PIGD          & \textbf{0.0000} & 0.0605 & --- \\
\bottomrule
\end{tabular}
\caption{Estimated rates per year for the true, unobserved phenotype.}
\end{table}

Going from TD straight to PIGD has a rate of 0.0057 per year. Going from TD to
Indeterminate has a rate of 0.0937 --- \textbf{sixteen times higher}. And going
from PIGD directly back to TD is estimated at exactly zero.

\textbf{Why this matters.} We never told the model the three states are ordered.
It worked that out on its own. And it is what you would expect if the underlying
thing is a continuous ratio: to get from above 1.15 to below 0.90, you have to
pass through the range in between.

\subsection*{The clinical label is unreliable, and we can say where}

\begin{table}[H]
\centering
\begin{tabular}{@{}lrrr@{}}
\toprule
\textbf{If the true state is\ldots} & \textbf{recorded TD} & \textbf{recorded Indet.} & \textbf{recorded PIGD} \\
\midrule
TD            & \textbf{94.1\%} & 3.8\% & 2.1\% \\
Indeterminate & 25.0\% & \textbf{50.0\%} & 25.0\% \\
PIGD          & 10.6\% & 5.2\% & \textbf{84.2\%} \\
\bottomrule
\end{tabular}
\end{table}

TD and PIGD are recorded correctly most of the time. Indeterminate is recorded
correctly at most half the time.

\textbf{Why this is convincing.} The previous fortnight found, by a completely
different method --- counting how often a change reverses --- that changes
\emph{into} Indeterminate reverse most often, 50 to 56\% of the time. Two
independent analyses landed on the same conclusion about the same state.

\subsection*{And a consequence: the phenotype is far more stable than its label}

\begin{table}[H]
\centering
\begin{tabular}{@{}lrr@{}}
\toprule
\textbf{State} & \textbf{Looking at the label} & \textbf{The true phenotype} \\
\midrule
TD            & 3.4 years & \textbf{10.1 years} \\
Indeterminate & 0.4 years & \phantom{0}\textbf{2.1 years} \\
PIGD          & 2.3 years & \textbf{16.5 years} \\
\bottomrule
\end{tabular}
\caption{Average time before leaving a state. Treating the label as the truth
makes the disease look like it changes three to seven times faster than it
does.}
\end{table}

% =============================================================================
\section{Step 6 --- We implemented the proposal's Bayesian model}
% =============================================================================

\textbf{What we did.} The hierarchical Bayesian growth model written out on
pages 13--15 of the proposal, exactly as specified. Each patient gets their own
starting level and their own rate of change, both drawn from a population
distribution. We wrote the sampler ourselves --- Gibbs sampling, about a second
to run --- and checked it by generating data from known values and confirming it
recovers them.

\textbf{Why.} It was in the plan, and it answers a different question. The
Markov model says \emph{whether} a patient will change. This one says \emph{how
fast their ratio is moving}, on the clinical scale, with a confidence interval.

\textbf{What we found.}

\begin{found}
\textbf{The TD/PIGD ratio falls by 0.093 per year} (95\% credible interval 0.068
to 0.116). The probability that the population is declining is essentially 1.

Starting from the average baseline ratio of 2.29, the average patient reaches
the TD threshold after about \textbf{12.3 years}. Patients differ: the
between-patient spread in rate is 0.103 per year, about as large as the average
rate itself.
\end{found}

\begin{figure}[H]
\centering
\includegraphics[width=\textwidth]{growth_trajectories.png}
\caption{Left: the fitted population trajectory, with 60 individual patients
drawn from the model. Right: the posterior for the rate of decline.}
\end{figure}

% =============================================================================
\section{Step 7 --- We found two problems in the specification}
% =============================================================================

Both were found by implementing the model, not by reading it. Present them as
contributions, not complaints.

\subsection*{Problem 1: two of the model's terms cannot be told apart}

The proposal has two separate error terms --- measurement error and process
noise. If you substitute one equation into the other, both are added to the same
observation, so the data can only ever see their \emph{sum}. We confirmed this
numerically: hold the total fixed, move any amount between the two, and the
likelihood does not change at all.

\textbf{Why it matters.} The headline output $P(\mathrm{TD}) = P(R^{true} >
1.15)$ refers to the \emph{true} ratio, whose spread excludes measurement error.
Without being able to separate the two, that quantity is not identified.

\textbf{What we did about it.} We estimate and report only the total, and say so
explicitly, rather than reporting a split the data cannot support.

\subsection*{Problem 2: a normal distribution does not fit a ratio}

The TD/PIGD ratio cannot be negative, is strongly skewed to the right, and is
exactly zero at 7\% of visits. A symmetric normal distribution cannot represent
that. The consequence is measurable: \textbf{12.2\% of the model's five-year
predictions are for a negative ratio}, which is impossible.

\ifasked{Are you saying the proposal's model is wrong?}{No --- it is a
reasonable model and it gives the interpretable progression rate that nothing
else gives. These are two specification details that only show up when you fit
it, and both have clean fixes: model the ratio on a log-like scale, and use the
state-space form the proposal also mentions, in which the two error terms
\emph{do} become separable.}

% =============================================================================
\section{Step 8 --- We compared the two models fairly}
% =============================================================================

\textbf{What we did.} Fitted both on the same training patients, and scored both
on exactly the same held-out transitions --- only those where every visit
involved has a usable ratio, so that neither model is given transitions the
other cannot attempt.

\textbf{What we found.} The hidden Markov model predicts subtype better, in
\textbf{6 splits out of 6}, by $+0.058 \pm 0.019$ (paired $t$-test $p =
6.6\times10^{-4}$).

\textbf{And we can say exactly why.}

\begin{found}
The growth model's residual spread is $\sigma = 1.138$. The Indeterminate band
is $1.15 - 0.90 = 0.25$ wide.

\textbf{The noise is 4.6 times wider than the band the model is being asked to
place a patient inside.} Even a perfectly estimated trajectory would land the
next observation almost anywhere.
\end{found}

\textbf{The conclusion we drew.} They are not competitors. The Markov model is
the better classifier; the growth model is the better description of
progression, and it is the one that leads to the second half of the project.
Keep both.

% =============================================================================
\section{Step 9 --- We realised we had been testing the wrong thing}
% =============================================================================

\textbf{What prompted it.} The clinical purpose is to help a doctor with a
patient in front of them --- often a patient being seen for the \emph{first}
time.

Every model above earns its accuracy by following a patient's history forward. A
new patient has no history. We had never tested that case.

\textbf{What we did.} Re-ran the comparison using only each patient's
\emph{first} visit, predicting their subtype at each later visit.

\textbf{What we found.}

\begin{table}[H]
\centering
\begin{tabular}{@{}lrr@{}}
\toprule
\textbf{Using only the first visit} & \textbf{Score} & \textbf{Accuracy} \\
\midrule
Baseline subtype $+$ ratio $+$ clinical measures & $\mathbf{-0.7550}$ & \textbf{71.9\%} \\
Baseline subtype $+$ clinical measures & $-0.7625$ & \\
Baseline subtype $+$ ratio & $-0.7735$ & \\
Baseline subtype alone & $-0.7813$ & \\
Hidden Markov model & $-0.7869$ & 68.1\% \\
Assume no change & $-1.4114$ & 64.6\% \\
\bottomrule
\end{tabular}
\end{table}

\begin{found}
\textbf{On the question that actually matters clinically, our best model was the
worst of the five.} Not because it is bad, but because it had nothing to work
with.

What a new patient \emph{does} have is a full baseline assessment. A direct
model using that beats the hidden Markov model by $+0.032$, at 71.9\% accuracy,
and is well calibrated.
\end{found}

\textbf{The rule this gives us.} The hidden Markov model's advantage grows as
follow-up accumulates: $+0.005$ after one visit, $+0.021$ after two, $+0.047$
after six or more. So: use the baseline model at the first visit, and switch to
the hidden Markov model once the patient has been seen a few times.

\ifasked{So was the Markov work wasted?}{No. It is the better model for every
visit after the first, which is most of the visits in a patient's care. What
changed is that we now know which model to use when, and we know it from
measurement rather than assumption.}

% =============================================================================
\section{Step 10 --- What a doctor would actually be shown}
% =============================================================================

\begin{table}[H]
\centering
\begin{tabular}{@{}lrrr@{}}
\toprule
\textbf{Horizon} & \textbf{P(TD)} & \textbf{P(Indeterminate)} & \textbf{P(PIGD)} \\
\midrule
now      & 100\% & 0\% & 0\% \\
1 year   & 82.1\% & 9.4\% & 8.6\% \\
2 years  & 78.5\% & 10.5\% & 11.0\% \\
3 years  & 74.9\% & 11.3\% & 13.8\% \\
5 years  & 67.1\% & 12.2\% & 20.7\% \\
\bottomrule
\end{tabular}
\caption{A new TD patient with an otherwise typical baseline profile. This is
the output format the proposal asks for.}
\end{table}

\noindent
The model is well calibrated, which is what makes the numbers usable:

\begin{table}[H]
\centering
\small
\begin{tabular}{@{}lrrrrr@{}}
\toprule
Predicted chance of PIGD & 4.5\% & 10.7\% & 21.4\% & 40.9\% & 66.8\% \\
What actually happened   & 2.0\% & 14.0\% & 24.5\% & 42.5\% & 63.7\% \\
\bottomrule
\end{tabular}
\end{table}

% =============================================================================
\section{Step 11 --- One result that runs against the plan}
% =============================================================================

We tested which baseline measures actually help, by removing each one and seeing
how much the held-out prediction worsens.

\begin{table}[H]
\centering
\small
\begin{tabular}{@{}lr|lr@{}}
\toprule
\textbf{Measure} & \textbf{Value} & \textbf{Measure} & \textbf{Value} \\
\midrule
MDS-UPDRS Part II & $+0.0153$ & MoCA (cognition) & $-0.0005$ \\
DaTscan putamen   & $+0.0050$ & UPSIT (smell)    & $-0.0008$ \\
Age               & $+0.0042$ & Disease duration & $-0.0010$ \\
Hoehn \& Yahr     & $+0.0021$ & SCOPA-AUT (autonomic) & $-0.0014$ \\
Sex               & $+0.0002$ & RBDSQ (sleep disorder) & $-0.0016$ \\
\bottomrule
\end{tabular}
\caption{Positive means the measure is earning its place.}
\end{table}

\begin{found}
\textbf{The four non-motor markers the second half of the project is built on
--- RBDSQ, SCOPA-AUT, UPSIT and MoCA --- all sit at or below zero.} The useful
information came from the motor and imaging measures instead.
\end{found}

\textbf{How to present this honestly.} It is not a refutation. The hypothesis in
the proposal is that these markers predict trajectory \emph{stability}
($\sigma_i$), and this test is about the trajectory \emph{mean}. Also, the
measures are correlated, so a genuinely useful marker can look worthless if
another measure already carries the same information.

But it is worth knowing before we build on them.

% =============================================================================
\section{Step 12 --- We finished the remaining pieces of the plan}
% =============================================================================

\begin{itemize}
  \item \textbf{Exploratory analysis of the ratio.} Done properly this time on
        the continuous ratio rather than the label. It shows the noise problem
        before any model is fitted --- within-patient spread 1.22 against a band
        0.25 wide.

  \item \textbf{Per-patient results.} We now produce, for every patient at every
        visit, the estimated true phenotype with a credible interval, plus each
        patient's own starting level and rate of change.

  \item \textbf{Validation by hiding the future.} Train on each patient's first
        two visits, predict all their later ones: 64.7\% accuracy over 1{,}886
        future visits, and 92.9\% of observations fell inside the 95\%
        prediction interval --- so the model's stated uncertainty is honest.
\end{itemize}

\subsection*{An important caveat we should raise ourselves}

\begin{found}
We measured how much of each patient's ``personal'' progression rate really
comes from their own data, rather than being the group average.

The answer is \textbf{5\%} for the typical patient. \textbf{No patient in the
cohort exceeds 80\%.} Their \emph{starting level}, by contrast, is well
determined for essentially everyone.

So the honest claim is that a prediction is personalised \textbf{by where the
patient is, not by how fast they are moving.} That is still worth a lot --- it
accounts for 95\% of the benefit of personalising at all --- but it is a weaker
claim than ``this patient progresses faster than average'', and we should not
make the stronger one.
\end{found}

\textbf{Why we are raising it.} The second half of the project wants to estimate
a patient-specific \emph{variability}. That needs more information than a rate
does, and the typical patient has only three visits. It would need the 182
patients with six or more visits, and that is better settled now than after
building on it.

% =============================================================================
\section{How to close}
% =============================================================================

\begin{enumerate}
  \item We built the model that was asked for. It works, and it is the baseline
        everything else had to beat.

  \item The obvious improvement failed. We tested why instead of guessing, and
        the test showed the observed subtype is a noisy reading of a hidden
        state --- $p = 7.9\times10^{-33}$.

  \item The model built on that finding wins 10 out of 10 validation splits, and
        it recovers the ordering of the three subtypes without being told it.

  \item We implemented the proposal's Bayesian model, got the progression rate
        with a confidence interval, and found two specification issues with
        clean fixes.

  \item We then realised the clinical question is a \emph{new} patient, tested
        that specifically, and found a different model wins there --- so we know
        which to use when.

  \item Every number is on patients the model never saw, repeated across
        independent splits, and reproducible from a single command. 144 automated
        tests cover the code.
\end{enumerate}

\vspace{0.8em}
\begin{found}
\textbf{The three questions most likely to come up.}

\smallskip
\emph{``Is the complicated model just overfitting?''} --- Every result is on
held-out patients, and the win repeats on ten independent splits.

\smallskip
\emph{``You said the interval matters, then the interval model failed. Which is it?''}
--- Both. The interval does matter; correcting for it inside a model that is
wrong in a more basic way makes things worse, and Step 3 identifies that basic
problem.

\smallskip
\emph{``Does a subtype change mean the disease actually progressed?''} --- Not on
its own; a third of observed changes reverse immediately. But the movement is
significantly one-directional ($p = 1.8\times10^{-8}$), so there is a real drift
underneath the noise. That is the case for this whole modelling approach.
\end{found}

\end{document}
